% =============================================================================
\section{Sistemas lineales e invariantes en el tiempo}
% =============================================================================

\begin{frame}{Propiedades de los sistemas}
  Un sistema $\mathcal{T}$ transforma una entrada en una salida: $y(t) = \mathcal{T}\{x(t)\}$.

  \begin{definicion}[linealidad]
    $\mathcal{T}\{a\,x_1(t) + b\,x_2(t)\} = a\,\mathcal{T}\{x_1(t)\} + b\,\mathcal{T}\{x_2(t)\}$
    para todo $a, b \in \C$.
  \end{definicion}

  \begin{definicion}[invariancia en el tiempo]
    Si $y(t) = \mathcal{T}\{x(t)\}$, entonces $y(t - t_0) = \mathcal{T}\{x(t - t_0)\}$ para todo $t_0$.
  \end{definicion}

  \begin{importante}
    Un sistema LTI queda completamente caracterizado por su
    \alert{respuesta al impulso} $h(t) = \mathcal{T}\{\impulso{t}\}$.
  \end{importante}
\end{frame}

\begin{frame}{Integral de convolución}
  \begin{columns}[c]
    \begin{column}{0.5\textwidth}
      \begin{teorema}[convolución]
        \[
          y(t) = x(t) \conv h(t) = \int_{-\infty}^{\infty} x(\tau)\, h(t - \tau) \dd \tau
        \]
      \end{teorema}
      \begin{itemize}
        \item<2-> Reflejar $h(\tau) \to h(-\tau)$
        \item<3-> Desplazar $t$: $h(t - \tau)$
        \item<4-> Integrar el producto (área sombreada)
      \end{itemize}
    \end{column}
    \begin{column}{0.46\textwidth}
      \centering
      \begin{tikzpicture}
        \begin{axis}[ejes senal, width=\linewidth, height=4.6cm,
            xlabel={$\tau$}, xmin=-1.5, xmax=3, ymin=0, ymax=1.3,
            xtick={0,2}, extra x ticks={1.2}, extra x tick labels={$t$},
            extra x tick style={tick label style={vision}}, ytick={1}]
          % x(τ): pulso rectangular en [0, 2]
          \addplot[senal continua, const plot] coordinates {(-1.5,0) (0,1) (2,0) (3,0)};
          % h(t-τ) con h(t) = e^{-t} u(t), para t = 1,2
          \only<3->{\addplot[energia, very thick, samples=100, domain=-1.5:1.2] {exp(-(1.2-x))};
                    \addplot[energia, very thick, domain=1.2:3] {0};}
          \only<4->{\addplot[fill=celeste, fill opacity=0.35, draw=none, domain=0:1.2, samples=60]
                      {exp(-(1.2-x))} \closedcycle;}
          \node[azulutn, anchor=south, font=\small] at (axis cs:1.6,1) {$x(\tau)$};
          \only<3->{\node[energia!85!black, anchor=south east, font=\small]
                      at (axis cs:1.15,1.02) {$h(t-\tau)$};}
        \end{axis}
      \end{tikzpicture}
    \end{column}
  \end{columns}
\end{frame}

\begin{frame}{Suma de convolución}
  Para sistemas LTI de tiempo discreto:
  \[
    y[n] = x[n] \conv h[n] = \sum_{k=-\infty}^{\infty} x[k]\, h[n - k]
  \]
  \begin{columns}[c]
    \begin{column}{0.5\textwidth}
      \begin{ejemplo}
        Con $x[n] = \escalon{n} - \escalon{n-4}$ y $h[n] = 0{,}5^{n}\escalon{n}$:
        \[
          y[n] = \sum_{k=0}^{\min(n,3)} 0{,}5^{\,n-k}, \quad n \geq 0
        \]
      \end{ejemplo}
    \end{column}
    \begin{column}{0.46\textwidth}
      \centering
      \begin{tikzpicture}
        \begin{axis}[ejes senal, width=\linewidth, height=4cm,
            xlabel={$n$}, ylabel={$y[n]$}, xmin=-2, xmax=10, ymin=0, ymax=2,
            xtick={0,2,...,10}, ytick={1, 1.875}, yticklabels={$1$, $1{,}875$}]
          \addplot[senal discreta] coordinates {
            (-2,0) (-1,0) (0,1) (1,1.5) (2,1.75) (3,1.875) (4,0.9375)
            (5,0.46875) (6,0.234375) (7,0.1171875) (8,0.05859375) (9,0.029296875) (10,0.0146484375)};
        \end{axis}
      \end{tikzpicture}
    \end{column}
  \end{columns}
\end{frame}

\begin{frame}{Interconexión de sistemas LTI}
  \centering
  \begin{tikzpicture}[node distance=1.1cm and 1.3cm]
    % Serie
    \node (x1) {$x(t)$};
    \node[bloque, right=of x1] (h1) {$h_1(t)$};
    \node[bloque, right=of h1] (h2) {$h_2(t)$};
    \node[right=of h2] (y1) {$y(t)$};
    \draw[flecha] (x1) -- (h1);
    \draw[flecha] (h1) -- (h2);
    \draw[flecha] (h2) -- (y1);
    \node[right=0.6cm of y1, anchor=west] {$h(t) = h_1(t) \conv h_2(t)$};

    % Paralelo
    \node[below=2.3cm of x1] (x2) {$x(t)$};
    \node[rama, right=0.7cm of x2] (r) {};
    \node[bloque, above right=0.25cm and 0.9cm of r] (g1) {$h_1(t)$};
    \node[bloque, below right=0.25cm and 0.9cm of r] (g2) {$h_2(t)$};
    \node[sumador, right=3.2cm of r] (s) {$+$};
    \node[right=0.9cm of s] (y2) {$y(t)$};
    \draw[thick] (x2) -- (r);
    \draw[flecha] (r) |- (g1);
    \draw[flecha] (r) |- (g2);
    \draw[flecha] (g1) -| (s);
    \draw[flecha] (g2) -| (s);
    \draw[flecha] (s) -- (y2);
    \node[right=0.6cm of y2, anchor=west] {$h(t) = h_1(t) + h_2(t)$};
  \end{tikzpicture}
\end{frame}
